\documentclass[11pt,a4paper]{article}

\usepackage[margin=1in]{geometry}
\usepackage{amsmath,amssymb}
\usepackage{booktabs}
\usepackage{graphicx}
\usepackage{hyperref}
\usepackage{enumitem}
\usepackage{siunitx}
\usepackage{caption}

\title{%
  \textbf{Floor Material Impact on Table-Tennis Ball Bounce Dynamics} \\[6pt]
  \large A Low-Order ODE Simulation of Topspin Bounce on Four Floor Types
}
\author{TableTennisSims Project}
\date{\today}

\begin{document}
\maketitle

\begin{abstract}
This document describes the mathematical model and simulation implementation
used to study how floor material beneath a table-tennis table affects ball
bounce dynamics. The simulation focuses on a \textbf{topspin bounce scenario}
and tracks four key observables across four floor types (concrete, hardwood,
rubber mat, carpet): (1)~horizontal ball velocity, (2)~vertical ball position
(post-bounce height), (3)~spin evolution, and (4)~vertical displacement of
the table/floor spring system. The model uses four degrees of freedom---three
for the ball and one for the table---with all contact interactions formulated
as smooth, continuous forces suitable for adaptive ODE solvers.
\end{abstract}

\tableofcontents
\newpage

%=============================================================================
\section{Introduction and Objectives}
%=============================================================================

The goal of this simulation is to predict how \textbf{floor material} beneath a
table-tennis table alters the outcome of a topspin ball bounce. The study
focuses on four primary observables:

\begin{enumerate}[nosep]
  \item \textbf{Horizontal velocity ($\dot{x}$):} How the ball's horizontal
    speed changes during contact, driven by the coupling between topspin and
    friction.
  \item \textbf{Ball height ($y$):} The vertical position of the ball over
    time, specifically how the post-bounce apex height varies across floor
    types.
  \item \textbf{Spin ($\omega$):} How the ball's topspin is altered by the
    bounce---friction at the contact patch converts between translational and
    rotational kinetic energy.
  \item \textbf{Table/floor displacement ($y_t$):} The vertical displacement
    of the table mass on the floor spring--dashpot, which is the primary
    quantity that differentiates floor materials in the simulation.
\end{enumerate}

The simulation scenario is a ball released from \SI{30.5}{cm} with an initial
horizontal velocity of \SI{2.0}{m/s} and topspin of \SI{150}{rad/s},
bouncing on a Butterfly Centrefold~25 tournament table (\SI{127}{kg}).

The model is intentionally \emph{low-order}: it retains only the degrees of
freedom and force laws necessary to capture the dominant physics, while
remaining numerically tractable with standard ODE integration techniques.

%=============================================================================
\section{Physical System and Assumptions}
%=============================================================================

\subsection{Ball}

The ball is modeled as a \textbf{rigid sphere} of mass~$m$, radius~$R$, and
moment of inertia~$I$. For a thin-walled hollow sphere,
\begin{equation}
  I = \frac{2}{3}\,m R^2.
  \label{eq:moi}
\end{equation}
The ball possesses translational degrees of freedom in the horizontal~($x$) and
vertical~($y$) directions, and a rotational degree of freedom~($\omega$)
about the axis perpendicular to the plane of motion (i.e.\ topspin/backspin
axis).

\subsection{Ball--Table Contact}

Deformation during contact (both of the ball shell and the very local table
surface) is represented by a \textbf{linear spring--dashpot} (Kelvin--Voigt
element) acting in the normal (vertical) direction. This element is active only
when geometric overlap exists. Tangential interaction at the contact patch is
governed by \textbf{Coulomb friction} with a smooth regularisation.

\subsection{Table and Legs}

The table top and its legs are lumped into a \textbf{single rigid body of
finite mass}~$m_t = \SI{127}{kg}$ (based on a Butterfly Centrefold~25
tournament table) with a single vertical translational degree of
freedom~$y_t$. No bending, torsion, or spatial wave propagation in the table
top is modeled; the table serves solely as a force-transmission path between
the ball and the floor support.

\subsection{Floor / Support}

The floor is modeled as a \textbf{second linear spring--dashpot} that acts on
the table mass. This element is always active (the table always rests on the
floor). The stiffness~$k_f$ and damping coefficient~$c_f$ of this element are
the \textbf{primary parameters that distinguish one floor type from another}.

\subsection{What Is Intentionally Not Modeled}

\begin{itemize}[nosep]
  \item Table-top bending or vibrational modes,
  \item Elastic wave propagation through the table or floor,
  \item Spatial extent or shape of the contact patch,
  \item Detailed viscoelastic or hyperelastic material constitutive laws,
  \item Spin decay during free flight,
  \item Aerodynamic drag and Magnus effect (identical across floor types).
\end{itemize}

These effects contribute negligibly to centimetre-level bounce accuracy and
post-bounce velocity/spin outcomes for the scenarios considered.

%=============================================================================
\section{Degrees of Freedom}
%=============================================================================

The system has \textbf{four degrees of freedom}:

\begin{center}
\begin{tabular}{@{}cll@{}}
  \toprule
  DOF & Symbol & Description \\
  \midrule
  1 & $x$     & Ball horizontal position \\
  2 & $y$     & Ball vertical position (measured from floor) \\
  3 & $\omega$ & Ball angular velocity (positive = topspin) \\
  4 & $y_t$   & Table vertical displacement from equilibrium \\
  \bottomrule
\end{tabular}
\end{center}

The associated velocity-level state variables are $\dot{x}$, $\dot{y}$,
$\dot{\omega}$ (angular acceleration enters through torque), and $\dot{y}_t$.

The full state vector used for numerical integration is
\begin{equation}
  \mathbf{s} =
  \begin{bmatrix}
    x \\ \dot{x} \\ y \\ \dot{y} \\ \omega \\ y_t \\ \dot{y}_t
  \end{bmatrix}
  \in \mathbb{R}^7.
  \label{eq:state}
\end{equation}

%=============================================================================
\section{Contact Geometry}
%=============================================================================

Let the table surface be located at vertical coordinate~$y_t$ (measured as a
displacement from its static equilibrium position, which we take as the
reference). The ball centre is at height~$y$. Contact occurs when the ball
surface geometrically overlaps the table surface:

\begin{equation}
  \delta = R - (y - y_t).
  \label{eq:overlap}
\end{equation}

Contact is active when $\delta > 0$. The overlap rate is
\begin{equation}
  \dot{\delta} = -(\dot{y} - \dot{y}_t).
  \label{eq:overlap_rate}
\end{equation}

%=============================================================================
\section{Forces During Contact}
%=============================================================================

\subsection{Normal (Vertical) Contact Force}

The normal contact force follows a linear spring--dashpot law:
\begin{equation}
  F_n =
  \begin{cases}
    k_c\,\delta + c_c\,\dot{\delta}, & \text{if } \delta > 0 \text{ and }
      k_c\,\delta + c_c\,\dot{\delta} > 0, \\
    0, & \text{otherwise}.
  \end{cases}
  \label{eq:Fn}
\end{equation}

The constraint $F_n \geq 0$ prevents the contact model from generating
artificial adhesion (tensile force) during the separation phase of the
collision. The parameters are:
\begin{itemize}[nosep]
  \item $k_c$: contact stiffness (\si{N/m}), representing ball shell and local
    table surface compliance,
  \item $c_c$: contact damping coefficient (\si{N\,s/m}), representing energy
    dissipation during deformation.
\end{itemize}

\subsection{Tangential (Frictional) Contact Force}

\subsubsection{Slip Velocity}

The velocity of the ball's contact point relative to the table surface is
\begin{equation}
  v_{\text{slip}} = \dot{x} - \omega\,R.
  \label{eq:vslip}
\end{equation}

When $v_{\text{slip}} = 0$ the ball is in a pure rolling state at the contact
point. A nonzero $v_{\text{slip}}$ indicates sliding.

\subsubsection{Friction Force}

The tangential friction force is modeled using Coulomb's law with a smooth
$\tanh$ regularisation to avoid a discontinuity at $v_{\text{slip}} = 0$:
\begin{equation}
  F_t = -\mu\,F_n\;\tanh\!\left(\frac{v_{\text{slip}}}{v_0}\right),
  \label{eq:Ft}
\end{equation}
where:
\begin{itemize}[nosep]
  \item $\mu$: coefficient of kinetic friction between ball and table surface,
  \item $v_0$: regularisation velocity (\si{m/s}), a small positive constant
    controlling the sharpness of the friction transition (typically
    \SIrange{0.01}{0.1}{m/s}).
\end{itemize}

The $\tanh$ function smoothly transitions the friction force direction across
zero slip velocity, approximating the sign function while maintaining
continuous derivatives required for ODE solver stability.

\subsubsection{Effect on Translation and Spin}

The friction force simultaneously modifies horizontal translation and spin:
\begin{align}
  m\,\ddot{x}    &= F_t,       \label{eq:xddot} \\
  I\,\dot{\omega} &= -R\,F_t.   \label{eq:omegadot}
\end{align}

Note the opposite signs: friction that decelerates horizontal sliding
simultaneously accelerates spin (and vice versa). This coupling is the
mechanism by which topspin produces increased horizontal speed after a bounce
while losing angular velocity.

%=============================================================================
\section{Equations of Motion}
%=============================================================================

\subsection{Ball --- Vertical}

\begin{equation}
  m\,\ddot{y} = -m\,g + F_n,
  \label{eq:ball_vert}
\end{equation}

where $g$ is gravitational acceleration and $F_n$ is defined by
Eq.~\eqref{eq:Fn}. During free flight ($\delta \leq 0$), $F_n = 0$ and the
ball follows a parabolic trajectory.

\subsection{Ball --- Horizontal}

\begin{equation}
  m\,\ddot{x} = F_t,
  \label{eq:ball_horiz}
\end{equation}

where $F_t$ is defined by Eq.~\eqref{eq:Ft}. During free flight, $F_t = 0$
and horizontal velocity is constant.

\subsection{Ball --- Spin}

\begin{equation}
  I\,\dot{\omega} = -R\,F_t.
  \label{eq:ball_spin}
\end{equation}

During free flight, spin is constant (aerodynamic spin decay is neglected).

\subsection{Table --- Vertical}

\begin{equation}
  m_t\,\ddot{y}_t = F_n - \bigl(k_f\,y_t + c_f\,\dot{y}_t\bigr).
  \label{eq:table_vert}
\end{equation}

The table is driven downward by the contact force $F_n$ (Newton's third law
applied to the ball--table interaction) and restored by the floor
spring--dashpot. When no ball contact exists ($F_n = 0$), the table undergoes
free damped oscillation about its equilibrium position.

%=============================================================================
\section{Summary of the ODE System}
%=============================================================================

Combining all equations of motion with the state vector defined in
Eq.~\eqref{eq:state}, the system to be integrated is:

\begin{equation}
  \frac{d\mathbf{s}}{dt} =
  \begin{bmatrix}
    \dot{x} \\[4pt]
    F_t / m \\[4pt]
    \dot{y} \\[4pt]
    -g + F_n / m \\[4pt]
    -R\,F_t / I \\[4pt]
    \dot{y}_t \\[4pt]
    \bigl(F_n - k_f\,y_t - c_f\,\dot{y}_t\bigr) / m_t
  \end{bmatrix},
  \label{eq:ode_system}
\end{equation}

where $F_n$ and $F_t$ are evaluated from the current state via
Eqs.~\eqref{eq:Fn}--\eqref{eq:Ft}.

%=============================================================================
\section{Energy Considerations}
%=============================================================================

During contact, energy flows through the system as follows:

\begin{enumerate}[nosep]
  \item Ball kinetic energy is converted to elastic potential energy in the
    contact spring (compression of ball and table surface).
  \item The contact dashpot ($c_c$) dissipates energy irreversibly---this is the
    \textbf{primary loss mechanism} and governs the effective coefficient of
    restitution.
  \item The contact force simultaneously accelerates the table mass downward.
  \item The table compresses the floor spring, storing elastic energy.
  \item The floor dashpot ($c_f$) dissipates energy irreversibly---this is the
    \textbf{secondary, floor-dependent loss mechanism}.
  \item Stored elastic energy in both springs is returned to the ball and table
    during rebound.
\end{enumerate}

The total mechanical energy of the system,
\begin{equation}
  E = \underbrace{\tfrac{1}{2}m\dot{x}^2
      + \tfrac{1}{2}m\dot{y}^2
      + \tfrac{1}{2}I\omega^2
      + m\,g\,y}_{\text{ball}}
    + \underbrace{\tfrac{1}{2}m_t\dot{y}_t^2
      + \tfrac{1}{2}k_f y_t^2}_{\text{table + floor spring}}
    + \underbrace{\tfrac{1}{2}k_c\delta^2}_{\text{contact spring (if } \delta > 0\text{)}},
  \label{eq:energy}
\end{equation}
should be \textbf{monotonically non-increasing} in time. Any increase in total
energy during a simulation indicates a numerical or implementation error and
serves as a primary validation diagnostic.

%=============================================================================
\section{Model Parameters}
%=============================================================================

\subsection{Calibrated Values}

The contact parameters $k_c$ and $c_c$ were tuned via parameter sweep to
simultaneously match two ITTF targets for a vertical drop from \SI{30.5}{cm}
onto a rigid (concrete) floor:

\begin{itemize}[nosep]
  \item Coefficient of restitution (COR) $\approx 0.76$
    (achieved: $0.768$),
  \item Contact time $\approx \SI{1}{ms}$
    (achieved: $\SI{1.01}{ms}$).
\end{itemize}

\begin{table}[ht]
\centering
\caption{Ball parameters (ITTF regulation ball).}
\label{tab:ball}
\begin{tabular}{@{}llll@{}}
  \toprule
  Parameter & Symbol & Value & Unit \\
  \midrule
  Mass                & $m$ & $2.7 \times 10^{-3}$ & \si{kg} \\
  Radius              & $R$ & $20 \times 10^{-3}$  & \si{m} \\
  Moment of inertia   & $I$ & $\frac{2}{3}mR^2 \approx 7.2 \times 10^{-7}$
                                                     & \si{kg\,m^2} \\
  \bottomrule
\end{tabular}
\end{table}

\begin{table}[ht]
\centering
\caption{Contact parameters (ball--table interaction). Values of $k_c$ and
$c_c$ were calibrated to match ITTF COR and contact-time targets.}
\label{tab:contact}
\begin{tabular}{@{}llll@{}}
  \toprule
  Parameter & Symbol & Value & Unit \\
  \midrule
  Contact stiffness       & $k_c$ & $2.5 \times 10^4$  & \si{N/m} \\
  Contact damping         & $c_c$ & $1.6$              & \si{N\,s/m} \\
  Friction coefficient    & $\mu$ & $0.25$             & --- \\
  Regularisation velocity & $v_0$ & $0.01$             & \si{m/s} \\
  \bottomrule
\end{tabular}
\end{table}

\begin{table}[ht]
\centering
\caption{Table parameters (Butterfly Centrefold~25 tournament table).}
\label{tab:table}
\begin{tabular}{@{}llll@{}}
  \toprule
  Parameter & Symbol & Value & Unit \\
  \midrule
  Table mass (top + legs) & $m_t$ & $127$ & \si{kg} \\
  \bottomrule
\end{tabular}
\end{table}

\begin{table}[ht]
\centering
\caption{Floor support parameters---the experimental variable. Values
selected from the representative ranges for each material class.}
\label{tab:floor}
\begin{tabular}{@{}lllll@{}}
  \toprule
  Floor Type & $k_f$ (\si{N/m}) & $c_f$ (\si{N\,s/m}) & Character \\
  \midrule
  Concrete slab     & $5 \times 10^7$ & $5 \times 10^3$ & Very stiff, low loss \\
  Hardwood (sprung) & $5 \times 10^5$ & $5 \times 10^3$ & Moderate compliance \\
  Rubber sport mat  & $5 \times 10^4$ & $5 \times 10^3$ & Compliant, higher damping \\
  Carpet over pad   & $5 \times 10^3$ & $5 \times 10^2$ & Soft, high loss \\
  \bottomrule
\end{tabular}
\end{table}

\subsection{Simulation Scenario}

The primary simulation scenario uses the following initial conditions:

\begin{center}
\begin{tabular}{@{}lll@{}}
  \toprule
  Parameter & Value & Description \\
  \midrule
  $y_0$ & \SI{0.305}{m} & Initial ball height (ITTF drop height) \\
  $\dot{x}_0$ & \SI{2.0}{m/s} & Initial horizontal velocity \\
  $\omega_0$ & \SI{150}{rad/s} & Initial topspin \\
  $t_{\text{span}}$ & $0$--$\SI{0.5}{s}$ & Simulation duration \\
  \bottomrule
\end{tabular}
\end{center}

%=============================================================================
\section{Numerical Solution Strategy}
%=============================================================================

The ODE system in Eq.~\eqref{eq:ode_system} is integrated using SciPy's
\texttt{solve\_ivp} with the \texttt{Radau} (implicit Runge--Kutta) method.
An implicit method is preferred because the system is \textbf{transiently
stiff}: during the brief contact phase (${\sim}\SI{1}{ms}$), the large
contact stiffness~$k_c$ introduces fast dynamics that require implicit
treatment, while the flight phase is non-stiff and can be integrated with
large steps.

Solver settings:
\begin{itemize}[nosep]
  \item Relative tolerance: $\texttt{rtol} = 10^{-9}$,
  \item Absolute tolerance: $\texttt{atol} = 10^{-9}$,
  \item Maximum step size: $\texttt{max\_step} = \SI{1}{ms}$,
  \item Dense output enabled for smooth animation rendering.
\end{itemize}

%=============================================================================
\section{Simulation Outputs}
%=============================================================================

The simulation produces two visualisations:

\subsection{Floor Comparison (4-panel static plot)}

A $2 \times 2$ panel comparing all four floor types across the key
observables:

\begin{enumerate}[nosep]
  \item \textbf{Ball height vs.\ time} --- Shows the parabolic flight,
    contact dip, and post-bounce apex. Any floor-dependent change in bounce
    height would appear here.
  \item \textbf{Spin vs.\ time} --- Shows the step decrease in $\omega$
    during contact as topspin is converted to horizontal speed via friction.
  \item \textbf{Horizontal velocity vs.\ time} --- Shows the step increase in
    $\dot{x}$ during contact (topspin-to-speed conversion). This is the
    primary ball-level observable of interest.
  \item \textbf{Table displacement vs.\ time} --- Shows $y_t$ in
    micrometres. This is the quantity that most clearly differentiates floor
    types: carpet produces ${\sim}\SI{9}{\micro m}$ of table displacement,
    while concrete produces sub-micrometre motion.
\end{enumerate}

\subsection{Bounce Animation (GIF)}

A 2D animated trajectory of the topspin bounce on the default (concrete)
floor, rendered in slow motion ($0.1\times$ real time). The animation
includes:
\begin{itemize}[nosep]
  \item Ball position (circle patch) with trajectory trail,
  \item Spin indicator (rotating radius line),
  \item Velocity vector arrow,
  \item Table surface line,
  \item Time display.
\end{itemize}

%=============================================================================
\section{Key Findings}
%=============================================================================

\subsection{Floor Effect on Ball Dynamics}

With a tournament-weight table ($m_t = \SI{127}{kg}$), the floor material has
\textbf{negligible effect on ball-level observables}. The ball height, spin,
and horizontal velocity curves overlap identically across all four floor
types. This is physically correct: the ball-to-table mass ratio is
approximately $1{:}47{,}000$, so the table acts as an effectively infinite
mass during the ${\sim}\SI{1}{ms}$ contact.

\subsection{Floor Effect on Table Response}

The floor effect manifests entirely in the \textbf{table displacement}
$y_t$. After the bounce:
\begin{itemize}[nosep]
  \item Carpet: ${\sim}\SI{9}{\micro m}$ peak displacement, slow recovery,
  \item Rubber mat: ${\sim}\SI{1.5}{\micro m}$ peak displacement,
  \item Hardwood: ${\sim}\SI{1}{\micro m}$ with moderate ringing,
  \item Concrete: sub-$\SI{0.1}{\micro m}$, rapid high-frequency ringing.
\end{itemize}

\subsection{Topspin Effect}

With $\omega_0 = \SI{150}{rad/s}$ and $\dot{x}_0 = \SI{2.0}{m/s}$, the slip
velocity at contact is $v_{\text{slip}} = 2.0 - 150 \times 0.02 =
\SI{-1.0}{m/s}$ (negative, meaning friction acts forward on the ball).
After the bounce, horizontal velocity increases from \SI{2.0}{m/s} to
${\sim}\SI{2.4}{m/s}$ while spin decreases from \SI{150}{rad/s} to
${\sim}\SI{120}{rad/s}$, confirming the expected topspin--speed coupling.

%=============================================================================
\section{Validation Summary}
%=============================================================================

\begin{center}
\begin{tabular}{@{}llll@{}}
  \toprule
  Check & Target & Result & Status \\
  \midrule
  COR (vertical drop, concrete) & $\approx 0.76$ & $0.768$ & Pass \\
  Contact time & $\approx \SI{1}{ms}$ & $\SI{1.01}{ms}$ & Pass \\
  Energy non-increasing & Monotonic & $< 10^{-8}$\,J noise & Pass \\
  Topspin gains horizontal speed & Yes & Yes & Pass \\
  Floor differentiation (table $y_t$) & Visible & Clear separation & Pass \\
  \bottomrule
\end{tabular}
\end{center}

%=============================================================================
\section{Implementation}
%=============================================================================

The simulation is implemented in Python using NumPy and SciPy, with
Matplotlib for visualisation. The code is structured as:

\begin{verbatim}
simulations/floor_material/
    parameters.py   -- Physical constants, material presets, SimConfig
    physics.py      -- ODE right-hand side (derivatives) and energy
    runner.py       -- solve_ivp integration wrapper
    visualize.py    -- Floor comparison plot and bounce animation
\end{verbatim}

To run:
\begin{verbatim}
source ~/.venv/bin/activate
python -m simulations.floor_material.visualize
\end{verbatim}

Output files are written to \texttt{output/}:
\begin{itemize}[nosep]
  \item \texttt{floor\_comparison.png} --- 4-panel floor comparison,
  \item \texttt{bounce\_animation.gif} --- Topspin bounce animation.
\end{itemize}

%=============================================================================
\section{Future Extensions}
%=============================================================================

The model is designed to be extended incrementally:
\begin{itemize}[nosep]
  \item \textbf{Aerodynamic drag:} Velocity-dependent drag force during flight,
    proportional to $v^2$ with Reynolds-dependent drag coefficient.
  \item \textbf{Magnus effect:} Spin-induced lift/lateral force during flight,
    enabling realistic topspin dip and sidespin curve.
  \item \textbf{Three-dimensional spin:} Extension to sidespin and gyroscopic
    effects (requires promoting $\omega$ to a vector).
  \item \textbf{Nonlinear contact:} Hertzian contact law
    ($F_n \propto \delta^{3/2}$) for improved accuracy at high impact
    velocities.
  \item \textbf{Multiple bounces:} Extended time span to capture cumulative
    floor effects over successive bounces.
  \item \textbf{Lighter table configurations:} Portable or home tables with
    lower mass where floor effects become more pronounced.
\end{itemize}

\end{document}
